\documentclass[12pt,a4paper]{article}
\usepackage[UTF8]{ctex}
\usepackage{amsmath,amssymb,amsthm}
\usepackage{geometry}
\usepackage{graphicx}

\geometry{margin=2.5cm}

\title{速度控制中$\theta_d$的含义：为什么通常不是0}
\author{第11章补充说明}
\date{}

\begin{document}

\maketitle

\section{公式11.11回顾}

速度输入的运动控制基本公式：
\begin{equation}
\dot{\theta} = \dot{\theta}_d + K_p(\theta_d - \theta) \tag{11.11}
\end{equation}

其中：
\begin{itemize}
\item $\dot{\theta}$: 实际关节速度（控制输出）
\item $\dot{\theta}_d$: 期望关节速度
\item $\theta_d$: \textbf{期望关节位置}（这是关键！）
\item $\theta$: 实际关节位置
\item $K_p$: 正定比例增益矩阵
\end{itemize}

\section{关键误解：$\theta_d$通常不是0}

\subsection{常见误解}

很多人认为$\theta_d$应该总是0，这是不正确的。这种误解可能来自：

\begin{enumerate}
\item \textbf{混淆了设定点控制和轨迹跟踪}：
\begin{itemize}
\item 设定点控制：$\theta_d$是常数（可能是0，也可能是其他值）
\item 轨迹跟踪：$\theta_d(t)$是时间的函数，随时间变化
\end{itemize}

\item \textbf{只考虑了特殊情况}：
\begin{itemize}
\item 如果只考虑"回到零位"的任务，$\theta_d = 0$
\item 但大多数机器人任务需要跟踪时变轨迹
\end{itemize}
\end{enumerate}

\subsection{$\theta_d$的实际含义}

$\theta_d(t)$是\textbf{期望的关节位置轨迹}，它：

\begin{itemize}
\item \textbf{通常是时变的}：$\theta_d = \theta_d(t)$，表示期望位置随时间变化
\item \textbf{由任务决定}：可以是任意轨迹，不一定是0
\item \textbf{与$\dot{\theta}_d$相关}：$\dot{\theta}_d = \frac{d\theta_d}{dt}$
\end{itemize}

\section{两种控制场景}

\subsection{场景1：设定点控制（Setpoint Control）}

\textbf{特点}：期望位置是常数，期望速度为零。

\begin{equation}
\theta_d = \text{常数}, \quad \dot{\theta}_d = 0
\end{equation}

\textbf{控制律简化为}：
\begin{equation}
\dot{\theta} = K_p(\theta_d - \theta) = K_p \theta_e
\end{equation}

\textbf{例子}：
\begin{itemize}
\item 机器人回到零位：$\theta_d = 0$
\item 机器人移动到特定位置：$\theta_d = \pi/2$（90度）
\item 机器人保持当前位置：$\theta_d = \theta(0)$（初始位置）
\end{itemize}

\textbf{注意}：即使在这种情况下，$\theta_d$也不一定是0，而是我们想要达到的目标位置。

\subsection{场景2：轨迹跟踪（Trajectory Tracking）}

\textbf{特点}：期望位置随时间变化，期望速度非零。

\begin{equation}
\theta_d = \theta_d(t), \quad \dot{\theta}_d = \frac{d\theta_d}{dt} \neq 0
\end{equation}

\textbf{控制律}：
\begin{equation}
\dot{\theta} = \dot{\theta}_d + K_p(\theta_d - \theta)
\end{equation}

\textbf{例子}：

\begin{enumerate}
\item \textbf{匀速运动}：
\begin{align}
\theta_d(t) &= \theta_0 + v_d t \\
\dot{\theta}_d(t) &= v_d \quad \text{（常数）}
\end{align}
其中$v_d$是期望速度。

\item \textbf{正弦运动}：
\begin{align}
\theta_d(t) &= A \sin(\omega t) \\
\dot{\theta}_d(t) &= A\omega \cos(\omega t)
\end{align}
其中$A$是振幅，$\omega$是角频率。

\item \textbf{多项式轨迹}：
\begin{align}
\theta_d(t) &= a_0 + a_1 t + a_2 t^2 + a_3 t^3 \\
\dot{\theta}_d(t) &= a_1 + 2a_2 t + 3a_3 t^2
\end{align}
这是常见的三次多项式轨迹。

\item \textbf{点到点运动}：
\begin{align}
\theta_d(t) &= \begin{cases}
\theta_1 & t < t_1 \\
\theta_1 + \frac{\theta_2 - \theta_1}{t_2 - t_1}(t - t_1) & t_1 \leq t \leq t_2 \\
\theta_2 & t > t_2
\end{cases}
\end{align}
从位置$\theta_1$移动到$\theta_2$。
\end{enumerate}

\section{控制律的物理意义}

让我们详细分析公式(11.11)的两部分：

\begin{equation}
\dot{\theta} = \underbrace{\dot{\theta}_d}_{\text{前馈项}} + \underbrace{K_p(\theta_d - \theta)}_{\text{反馈项}}
\end{equation}

\subsection{前馈项：$\dot{\theta}_d$}

\begin{itemize}
\item \textbf{作用}：提供期望速度，使系统能够跟踪运动轨迹
\item \textbf{特点}：这是"开环"项，不依赖于实际位置
\item \textbf{必要性}：如果没有这一项，系统只能跟踪固定位置，无法跟踪运动轨迹
\end{itemize}

\subsection{反馈项：$K_p(\theta_d - \theta)$}

\begin{itemize}
\item \textbf{作用}：根据位置误差修正速度，消除跟踪误差
\item \textbf{特点}：这是"闭环"项，依赖于实际位置
\item \textbf{必要性}：即使有前馈项，由于干扰、建模误差等原因，仍需要反馈来消除误差
\end{itemize}

\section{具体例子说明}

\subsection{例子1：跟踪固定位置（$\theta_d = 0$的特殊情况）}

\textbf{任务}：让机器人回到零位并保持。

\textbf{参数}：
\begin{itemize}
\item $\theta_d = 0$（常数）
\item $\dot{\theta}_d = 0$（因为位置不变）
\item 初始位置：$\theta(0) = 0.5$ rad
\end{itemize}

\textbf{控制律}：
\begin{equation}
\dot{\theta} = 0 + K_p(0 - \theta) = -K_p \theta
\end{equation}

\textbf{行为}：
\begin{itemize}
\item 如果$\theta > 0$，则$\dot{\theta} < 0$，机器人向零位移动
\item 如果$\theta < 0$，则$\dot{\theta} > 0$，机器人向零位移动
\item 最终$\theta \to 0$
\end{itemize}

\textbf{注意}：这是$\theta_d = 0$的特殊情况，但$\theta_d$也可以是其他常数，比如$\theta_d = \pi/2$。

\subsection{例子2：跟踪运动轨迹（$\theta_d \neq 0$且时变）}

\textbf{任务}：让机器人以1 rad/s的速度匀速运动。

\textbf{参数}：
\begin{itemize}
\item $\theta_d(t) = t$（位置随时间线性增加）
\item $\dot{\theta}_d(t) = 1$（常数速度）
\item 初始位置：$\theta(0) = 0$
\end{itemize}

\textbf{控制律}：
\begin{equation}
\dot{\theta} = 1 + K_p(t - \theta)
\end{equation}

\textbf{行为分析}：

假设在某个时刻$t = 2$ s：
\begin{itemize}
\item 期望位置：$\theta_d(2) = 2$ rad
\item 实际位置：$\theta(2) = 1.8$ rad（有误差）
\item 位置误差：$\theta_e = 2 - 1.8 = 0.2$ rad
\item 期望速度：$\dot{\theta}_d = 1$ rad/s
\end{itemize}

控制律计算：
\begin{equation}
\dot{\theta} = 1 + K_p \times 0.2 = 1 + 0.2K_p
\end{equation}

如果$K_p = 10$：
\begin{equation}
\dot{\theta} = 1 + 2 = 3 \text{ rad/s}
\end{equation}

\textbf{解释}：
\begin{itemize}
\item 前馈项（1 rad/s）提供基本速度
\item 反馈项（2 rad/s）加速追赶，消除位置误差
\item 当误差减小时，反馈项也减小，最终$\dot{\theta} \to 1$ rad/s
\end{itemize}

\subsection{例子3：跟踪正弦轨迹}

\textbf{任务}：让机器人跟踪正弦运动。

\textbf{参数}：
\begin{itemize}
\item $\theta_d(t) = 0.5 \sin(2t)$（振幅0.5 rad，频率1 rad/s）
\item $\dot{\theta}_d(t) = \cos(2t)$
\item 初始位置：$\theta(0) = 0$
\end{itemize}

\textbf{控制律}：
\begin{equation}
\dot{\theta} = \cos(2t) + K_p(0.5\sin(2t) - \theta)
\end{equation}

\textbf{行为}：
\begin{itemize}
\item 前馈项$\cos(2t)$提供期望速度
\item 反馈项根据位置误差修正速度
\item 系统跟踪正弦轨迹
\end{itemize}

\section{为什么需要$\theta_d$项？}

\subsection{如果没有$\theta_d$项会怎样？}

如果控制律只有：
\begin{equation}
\dot{\theta} = \dot{\theta}_d
\end{equation}

\textbf{问题}：
\begin{itemize}
\item 这是纯前馈控制，没有反馈
\item 如果初始位置有误差，误差会一直存在
\item 如果存在干扰，无法消除误差
\item 系统无法收敛到期望轨迹
\end{itemize}

\subsection{如果没有$\dot{\theta}_d$项会怎样？}

如果控制律只有：
\begin{equation}
\dot{\theta} = K_p(\theta_d - \theta)
\end{equation}

\textbf{问题}：
\begin{itemize}
\item 如果$\theta_d$是常数，可以工作（设定点控制）
\item 但如果$\theta_d$在变化，系统会"滞后"
\item 例如，如果$\theta_d(t) = t$，系统会一直"追赶"，但永远追不上
\item 需要$\dot{\theta}_d$项来提供"预测"速度
\end{itemize}

\subsection{两个项缺一不可}

\begin{itemize}
\item \textbf{$\dot{\theta}_d$项}：使系统能够跟踪运动轨迹
\item \textbf{$K_p(\theta_d - \theta)$项}：消除跟踪误差，提高鲁棒性
\item \textbf{两者结合}：既能跟踪运动轨迹，又能消除误差
\end{itemize}

\section{总结}

\subsection{关键要点}

\begin{enumerate}
\item \textbf{$\theta_d$通常不是0}：
\begin{itemize}
\item 在设定点控制中，$\theta_d$是常数（可能是0，也可能是其他值）
\item 在轨迹跟踪中，$\theta_d(t)$是时变函数
\end{itemize}

\item \textbf{$\theta_d$由任务决定}：
\begin{itemize}
\item 不同的任务有不同的$\theta_d(t)$
\item $\theta_d$可以是任意轨迹，不一定是0
\end{itemize}

\item \textbf{控制律的两部分}：
\begin{itemize}
\item $\dot{\theta}_d$：前馈项，提供期望速度
\item $K_p(\theta_d - \theta)$：反馈项，消除位置误差
\end{itemize}

\item \textbf{实际应用}：
\begin{itemize}
\item 大多数机器人任务需要跟踪时变轨迹
\item 因此$\theta_d(t)$通常是时间的函数，不是常数0
\end{itemize}
\end{enumerate}

\subsection{常见错误理解}

\begin{itemize}
\item \textbf{错误}：$\theta_d$总是0
\item \textbf{正确}：$\theta_d$由任务决定，可以是任意值或函数
\end{itemize}

\begin{itemize}
\item \textbf{错误}：速度控制只需要$\dot{\theta}_d$
\item \textbf{正确}：需要$\dot{\theta}_d$和$K_p(\theta_d - \theta)$两项
\end{itemize}

\begin{itemize}
\item \textbf{错误}：$\theta_d$和$\dot{\theta}_d$是独立的
\item \textbf{正确}：$\dot{\theta}_d = \frac{d\theta_d}{dt}$，两者相关
\end{itemize}

\section{数学关系}

\subsection{$\theta_d$和$\dot{\theta}_d$的关系}

如果$\theta_d(t)$是时间的函数，则：
\begin{equation}
\dot{\theta}_d(t) = \frac{d\theta_d(t)}{dt}
\end{equation}

\textbf{例子}：
\begin{align}
\text{如果} \quad \theta_d(t) &= t^2 \\
\text{则} \quad \dot{\theta}_d(t) &= 2t
\end{align}

\subsection{误差动力学}

定义位置误差：
\begin{equation}
\theta_e(t) = \theta_d(t) - \theta(t)
\end{equation}

对时间求导：
\begin{equation}
\dot{\theta}_e(t) = \dot{\theta}_d(t) - \dot{\theta}(t)
\end{equation}

代入控制律(11.11)：
\begin{align}
\dot{\theta}_e &= \dot{\theta}_d - [\dot{\theta}_d + K_p(\theta_d - \theta)] \\
              &= -K_p(\theta_d - \theta) \\
              &= -K_p \theta_e
\end{align}

\textbf{结论}：误差动力学是一阶系统：
\begin{equation}
\dot{\theta}_e + K_p \theta_e = 0
\end{equation}

解为：
\begin{equation}
\theta_e(t) = \theta_e(0) e^{-K_p t}
\end{equation}

这表明误差指数衰减到0，系统能够跟踪期望轨迹。

\end{document}

